\chapter{Module analytics et espace administrateur}

\section{Introduction}
Ce chapitre decrit la supervision fonctionnelle de la plateforme via les tableaux de bord et l'espace d'administration.

\section{Collecte des evenements analytics}
Le systeme enregistre les interactions importantes :
\begin{itemize}
  \item messages envoyes et reponses recues ;
  \item profils visiteurs (tranche d'age, nationalite, motif) ;
  \item actions d'administration.
\end{itemize}

\section{KPI et visualisations}
Les indicateurs permettent d'analyser :
\begin{itemize}
  \item usage du chatbot ;
  \item repartition demographique des utilisateurs ;
  \item activite par hotel et par periode ;
  \item performance des pages et interactions cles.
\end{itemize}

\section{Gestion administrative}
L'interface admin permet :
\begin{itemize}
  \item la gestion des parametres hoteliers ;
  \item la gestion des evenements et attractions ;
  \item le suivi des reservations ;
  \item la maintenance des contenus dynamiques.
\end{itemize}

\section{Authentification et controle d'acces}
L'acces administrateur est protege par JWT et verification serveur. Les routes sensibles sont controlees pour empecher les acces non autorises.

\section{Conclusion}
Ce module transforme les donnees d'usage en informations exploitables pour le pilotage operationnel. Le chapitre suivant traite des aspects securite, tests et deploiement.
